\chapter{Liver Disease}

\section*{Patient focus}
The liver is a vital organ that processes nutrients, filters blood, and produces bile. When the liver is damaged---by alcohol, viral hepatitis, or fatty deposits---imaging can show the extent of damage, detect scarring (cirrhosis), and find complications like fluid in the abdomen (ascites) or abnormal blood vessels. Ultrasound is often the first scan, used to check for fatty liver and screen for liver cancer. CT and MRI provide more detailed information when needed.


\label{ch:ch64_liver_disease}

\section{Imaging}
Imaging evaluates chronic liver disease, portal hypertension, focal lesions, vascular complications, and treatment response. Ultrasound, multiphasic CT, MRI, elastography, and Doppler are selected according to the question and risk group.

\section{Fatty Liver (NAFLD/NASH)}

\begin{itemize}
  \item \textbf{US findings}: increased hepatic echogenicity (``bright liver''), poor diaphragm visualization, loss of portal vein wall echoes
  \item \textbf{CT}: diffuse low attenuation of liver parenchyma (liver $<$ spleen attenuation on non-contrast CT); hepatic steatosis $>$33\%
  \item \textbf{MRI}: proton density fat fraction (PDFF) quantifies steatosis---gold standard for non-invasive fat quantification
  \item \textbf{NASH vs NAFLD}: imaging cannot reliably distinguish; NASH requires biopsy or advanced fibrosis markers
\end{itemize}

\section{Cirrhosis Imaging}

\subsection{CT and MRI Signs of Cirrhosis}
\begin{itemize}
  \item \textbf{Morphological}: nodular liver contour, hypertrophy of caudate and left lobes, atrophy of right lobe, widening of fissures
  \item \textbf{Parenchymal}: heterogeneous enhancement, regenerative nodules (iso- to hyperintense on T1), dysplastic nodules (variable), fibrosis (delayed enhancement)
  \item \textbf{Fibrosis on MRI}: T1 mapping and gadolinium-enhanced imaging can stage fibrosis
\end{itemize}

\begin{pearlbox}[Cirrhosis on imaging]
A nodular liver contour with caudate hypertrophy and right lobe atrophy is virtually diagnostic of cirrhosis. However, early cirrhosis may have a near-normal appearance on CT---ultrasound elastography or MRI is more sensitive for early fibrosis.
\end{pearlbox}

\section{Portal Hypertension}

Portal hypertension is defined as hepatic venous pressure gradient (HVPG) $\geq$5 mmHg; clinically significant at $\geq$10 mmHg.

\subsection{Imaging Signs}
\begin{itemize}
  \item \textbf{Direct}: portal vein diameter $>$13 mm, hepatofugal flow (flow away from liver on Doppler), portosystemic shunts
  \item \textbf{Indirect}: splenomegaly ($>$13 cm), ascites, varices (oesophageal, gastric, paraumbilical, splenorenal), gallbladder wall thickening
  \item \textbf{Recanalised paraumbilical vein}: pathognomonic of portal hypertension
\end{itemize}

\subsection{Baveno Classification}
\begin{table}[h]
\centering
\caption{Portal hypertension grading on imaging}
\begin{tabular}{|l|l|}
\hline
\textbf{Grade} & \textbf{Features} \\
\hline
Compensated & Varices, splenomegaly, no ascites \\
Decompensated & Ascites, variceal bleeding, hepatic encephalopathy \\
\hline
\end{tabular}
\end{table}

\section{Hepatocellular Carcinoma (HCC)---LI-RADS}

\begin{itemize}
  \item \textbf{Population}: cirrhotic patients, chronic hepatitis B, or other high-risk groups
  \item \textbf{Technique}: multiphasic CT or MRI with arterial phase hyperenhancement, portal venous/delayed washout, capsule appearance
  \item \textbf{LI-RADS categories}:
    \begin{enumerate}
      \item LR-1: definitely benign
      \item LR-2: probably benign
      \item LR-3: intermediate probability
      \item LR-4: probably HCC
      \item LR-5: definitely HCC (arterial hyperenhancement + washout + capsule + size $\geq$10 mm)
      \item LR-M: probably malignancy other than HCC
      \item LR-TIV: tumor in vein
    \end{enumerate}
\end{itemize}

\begin{pearlbox}[Screening for HCC]
All patients with cirrhosis should have 6-monthly ultrasound with or without AFP for HCC screening. If a lesion is detected, proceed to multiphasic CT or MRI and apply LI-RADS. LR-5 lesions in appropriate clinical context do not require biopsy for diagnosis.
\end{pearlbox}

\section{Liver Fibrosis Assessment---Elastography}

\subsection{Ultrasound Elastography}
\begin{itemize}
  \item \textbf{Transient elastography (FibroScan)}: measures liver stiffness in kPa; correlates with fibrosis stage
  \item \textbf{Shear wave elastography (SWE)}: uses acoustic radiation force; can be performed with standard US machines
  \item \textbf{Cutoffs (approximate)}: $<$7 kPa = F0-F1 (no/mild fibrosis); 7--10 kPa = F2; 10--14 kPa = F3; $>$14 kPa = F4 (cirrhosis)
\end{itemize}

\subsection{MRI Elastography}
\begin{itemize}
  \item More accurate than US elastography, especially in obesity and ascites
  \item Provides whole-liver stiffness map
  \item Cutoffs differ from US elastography
\end{itemize}

\section{Other Liver Diseases}

\begin{table}[h]
\centering
\caption{Imaging features of common liver diseases}
\begin{tabular}{|p{0.33\textwidth}|p{0.55\textwidth}|}
\hline
\textbf{Disease} & \textbf{Key imaging features} \\
\hline
Alcoholic liver disease & Fatty liver, cirrhosis, enlarged caudate \\
Autoimmune hepatitis & Cirrhosis, may have atrophic left lobe \\
Budd-Chiari & Hepatic vein occlusion, caudate hypertrophy, collateral vessels \\
Hemochromatosis & Diffuse low T1 and T2 signal (iron deposition) \\
Wilson's disease & Cirrhosis, copper deposition (variable imaging) \\
\hline
\end{tabular}
\end{table}

\section{Reporting}
Describe liver morphology, steatosis or iron, fibrosis indicators, portal and hepatic veins, collaterals, ascites, splenic size, and focal lesion enhancement. Apply LI-RADS only when the clinical population and technique meet its requirements.

\section{Safety}
Contrast choice and kidney function require current protocol review. Report portal-vein thrombosis, acute bleeding, biliary obstruction, abscess, or suspicious lesion promptly.

\begin{pearlbox}[Gadolinium safety in liver disease]
Many patients with cirrhosis have reduced renal function. If contrast-enhanced MRI is needed, use the current agent-specific ACR risk group and local protocol. Group II agents have a very low NSF risk and should not automatically be withheld solely for eGFR $<$30. Ultrasound-based techniques may answer some questions, but are not interchangeable with MRI for every indication.
\end{pearlbox}

\clearpage
